\chapter[HBEV]{Holographic Boundary-Entropy Vacuum and \(\Lambda_{\mathrm{eff}}\)}
\label{ch:hbev}

\paragraph{Learning outcomes.}
After this chapter the reader should be able to:
\begin{enumerate}
  \item restate the quartic vacuum-energy debt isolated in
        Chapter~8 as a problem about a bulk term
        \(\rho_\Lambda\sqrt{-g}\);
  \item write the HBEV boundary action built from the
        Gibbons--Hawking--York term and an entropy density
        \(S_{\mathrm{ent}}(h_{ij},\phi)\);
  \item obtain \(\Lambda_{\mathrm{eff}}\) from stationarity of that
        boundary action rather than from \(\LUV^4\);
  \item check that the regulated cosmological term still obeys
        Theorem~\ref{thm:conservation};
  \item enter Chapter~13 knowing which vacuum energy is being
        inserted into regular-core and galactic solutions.
\end{enumerate}

Chapter~8 produced
\begin{equation}
  \label{eq:rho-debt}
  \rho_\Lambda
  =-\sum_{s\in\mathcal{S}_{\mathrm{sec}}}
    \frac{w_s\operatorname{tr}_{V_s}\mathbb I_s}{64\pi^2}
    \bigl({\LUV}_s^4-\mu_s^4\bigr)+\cdots.
\end{equation}
This is the truncated, unresummed sector sum of
Chapter~8 under its stated continuation, with the same included
sectors and determinant weights. A common-limit form using
\(\mathcal N_0\) applies only when all included sectors are assigned
the same limits; otherwise sector-dependent thresholds are retained.
Only if the weighted quartic contribution is nonzero and remains
uncancelled, including after finite local terms and counterterms are
specified, must its magnitude be reconciled with the observed
cosmological constant. Without those conditions no such comparison
follows from~\eqref{eq:rho-debt}. This chapter does not cancel it by a
bare counterterm tuned to 120 decimal places. It replaces the bulk
cosmological term by a boundary variational problem whose solution
is an infrared scale.

\section{What would not count as a solution}
\label{sec:ch12-not}

Three common manoeuvres are excluded in advance.

\begin{enumerate}
  \item \textbf{Subtraction by hand.}
        Adding \(-\rho_\Lambda\int\sqrt{-g}\) to the action and
        declaring the sum zero is a renormalization condition at
        one scale. It does not explain the residual scale.
  \item \textbf{Setting \(\LUV\) to meV.}
        The cutoff constraint cannot be assessed until the physical
        \(G_{\mathrm{eff}}\) is obtained from the specified sector
        calculation of Problem~F.
  \item \textbf{A spacetime-dependent \(\rho_\Lambda(x)\) with no
        current.}
        Section~\ref{sec:ch11-vacuum} forbids this unless extra
        stress is supplied. Regulation that breaks
        Theorem~\ref{thm:conservation} is not available.
\end{enumerate}

Any construction that survives must therefore (i)~replace the bulk
quartic by an infrared datum, (ii)~use a \(G_{\mathrm{eff}}\) matched
from specified operators, weights, thresholds, continuation, and
scheme as required by Problem~F, and (iii)~produce a conserved
\(\Lambda_{\mathrm{eff}}g_{\mu\nu}\) or a conserved completion of
it.

\section{Boundary geometry}
\label{sec:ch12-boundary}

Let \(M\) be a region of the emergent spacetime with boundary
\(\partial M\), outward normal \(n^\mu\), induced metric
\begin{equation}
  \label{eq:induced}
  h_{ij}=g_{\mu\nu}e^\mu_i e^\nu_j,
\end{equation}
and extrinsic curvature
\begin{equation}
  \label{eq:K}
  K_{ij}=\nabla_{(\mu}n_{\nu)}e^\mu_i e^\nu_j,
  \qquad
  K=h^{ij}K_{ij}.
\end{equation}
The Gibbons--Hawking--York term already introduced
in~\eqref{eq:GHY} is
\begin{equation}
  \label{eq:GHY-again}
  S_{\mathrm{GHY}}
  =\frac{1}{8\pi\Geff}
   \int_{\partial M}\mathrm{d}^3x\,\sqrt{|h|}\,K.
\end{equation}
On a timelike boundary \(K\) measures the expansion of the outward
normal; on a spacelike boundary it measures the expansion of a
unit-normal congruence. HBEV uses a timelike conformal boundary, or
a large-radius timelike tube that will be sent to infinity after
variation, according to the problem.

The boundary may also carry condensate data. Write \(\phi\) for the
pull-back of the polar radial mode \(\sigma\) and of any unlocked
scalars that survive Assumption~\ref{ass:locked-block}. The pair
\((h_{ij},\phi)\) is the boundary field content.

\section{The HBEV action}
\label{sec:ch12-action}

\begin{definition}[HBEV boundary action]
  \label{def:HBEV}
  The holographic boundary-entropy vacuum functional is
  \begin{equation}
    \label{eq:SHBEV}
    S_{\mathrm{HBEV}}[h,\phi]
    =\int_{\partial M}\mathrm{d}^3x\,\sqrt{|h|}
     \Bigl(
       \frac{K}{8\pi\Geff}
       -T_{\mathrm{bdy}}
       -S_{\mathrm{ent}}(h_{ij},\phi)
     \Bigr),
  \end{equation}
  where \(T_{\mathrm{bdy}}\) is a boundary tension (a
  counterterm density built from \(h_{ij}\) only) and
  \(S_{\mathrm{ent}}(h_{ij},\phi)\) is an entanglement-entropy
  density of the condensate restricted to \(\partial M\).
\end{definition}

The first term is GHY. The second term is the standard counterterm
that renders the on-shell action finite on AdS-like or
asymptotically flat tubes; in the present setting it is fixed by
demanding that \(S_{\mathrm{HBEV}}\) vanish on the reference
condensate \(\bar e^A_\mu=\delta^A_\mu\), \(\phi=0\). The third term
is the new ingredient. It is not the diamond entropy \(S_D\) of
Chapter~6. It is a local density on \(\partial M\) whose integral
tracks the entanglement of bulk condensate degrees of freedom that
end on the boundary.

\begin{assumption}[Existence of a local entropy density]
  \label{ass:Sent-density}
  In the condensed phase there exists a diffeomorphism-scalar
  density \(S_{\mathrm{ent}}(h_{ij},\phi)\), of mass dimension
  three, obtained from the replica construction of
  Section~\ref{sec:ch6-replica} applied to a thin tubular
  neighbourhood of \(\partial M\) and then localized. Its leading
  terms are
  \begin{equation}
    \label{eq:Sent-exp}
    S_{\mathrm{ent}}
    =\alpha_0\LUV^3
     +\alpha_1\LUV R^{(h)}
     +\alpha_2\,\phi^2\LUV
     +\cdots,
  \end{equation}
  with dimensionless coefficients \(\alpha_i\) of order the fibre
  trace.
\end{assumption}

The \(\alpha_0\LUV^3\) piece is the boundary analogue of
\(\rho_\Lambda\). It will be cancelled against \(T_{\mathrm{bdy}}\)
and against the leading GHY contribution on the reference
background. What is not cancelled is required to be an infrared
quantity.

\section{Stationarity and the induced cosmological term}
\label{sec:ch12-stationarity}

Vary \(S_{\mathrm{EH}}+S_{\mathrm{HBEV}}\) with respect to the
boundary metric \(h_{ij}\), holding the bulk fields on shell. The
bulk Einstein tensor produces the usual Brown--York stress. The
boundary variation produces
\begin{equation}
  \label{eq:BY}
  \pi^{ij}
  =\frac{1}{8\pi\Geff}
   \bigl(K^{ij}-Kh^{ij}\bigr)
   -T_{\mathrm{bdy}}^{ij}
   -\frac{2}{\sqrt{|h|}}
    \frac{\delta S_{\mathrm{ent}}}{\delta h_{ij}}.
\end{equation}
HBEV imposes entanglement equilibrium at the boundary in the same
logical role that Assumption~\ref{ass:EE} played in the bulk
diamonds:

\begin{assumption}[Boundary entanglement equilibrium]
  \label{ass:bdy-EE}
  The physical vacuum is a stationary point of \(S_{\mathrm{HBEV}}\)
  at fixed boundary volume element, in the sense that the
  trace-adjusted boundary stress vanishes,
  \begin{equation}
    \label{eq:bdy-EE}
    h_{ij}\pi^{ij}=0.
  \end{equation}
\end{assumption}

On an asymptotically simple tube of radius \(r_B\), with
\(K=2/r_B+O(r_B^{-3})\) in asymptotically flat space and
\(K=3/\ell_{\mathrm{AdS}}+O(r_B^{-2})\) in AdS, equation
\eqref{eq:bdy-EE} reduces to a single condition on the integration
constant in the bulk Einstein equation with a cosmological term.
That constant is \(\Lambda_{\mathrm{eff}}\).

\begin{theorem}[Induced cosmological constant]
  \label{thm:Leff}
  Under Assumptions~\ref{ass:Sent-density} and~\ref{ass:bdy-EE},
  and with \(T_{\mathrm{bdy}}\) chosen to cancel the reference
  \(\alpha_0\LUV^3\) piece, the bulk field equation compatible
  with boundary stationarity is
  \begin{equation}
    \label{eq:EFE-Leff}
    G_{\mu\nu}+\Lambda_{\mathrm{eff}}g_{\mu\nu}
    =8\pi\Geff\,T_{\mu\nu}^{\mathrm{total}},
  \end{equation}
  where
  \begin{equation}
    \label{eq:Leff}
    \frac{\Lambda_{\mathrm{eff}}}{8\pi\Geff}
    =\bigl\langle
       S_{\mathrm{ent}}^{\mathrm{(IR)}}
     \bigr\rangle_{\partial M}
    -T_{\mathrm{bdy}}^{\mathrm{(IR)}}
  \end{equation}
  is determined by the infrared part of the boundary entropy
  density after the \(\LUV^3\) subtraction. In particular
  \(\Lambda_{\mathrm{eff}}\) does not scale as \(\LUV^2\).
\end{theorem}

\begin{proof}[Proof sketch]
  On-shell variation of \(S_{\mathrm{EH}}+S_{\mathrm{GHY}}\) with
  respect to a constant bulk cosmological term is proportional to
  the spacetime volume, which by Stokes is a boundary integral of
  the normal. That boundary integral is matched to
  \(h_{ij}\pi^{ij}\) after the GHY contribution is included.
  Imposing~\eqref{eq:bdy-EE} fixes the constant. The ultraviolet
  pieces of \(S_{\mathrm{ent}}\) are cancelled by the choice of
  \(T_{\mathrm{bdy}}\) on the reference background, which is a
  finite renormalization of a three-dimensional density, not a
  fine-tuning of a four-dimensional quartic against itself. The
  remainder is infrared by construction of the subtraction, and
  is constant on the vacuum because Assumption~\ref{ass:bdy-EE}
  is imposed uniformly on the asymptotic tube.
\end{proof}

\begin{remark}[What the theorem does not claim]
  Equation~\eqref{eq:Leff} does not compute the observed value
  \(10^{-122}\) in Planck units. It changes the \emph{origin} of
  the scale: from \(\LUV^4\) to an infrared boundary entropy.
  Predicting the numerical value requires an evaluation of
  \(S_{\mathrm{ent}}^{\mathrm{(IR)}}\) on the actual cosmological
  boundary, which this book does not perform. Chapter~17 lists
  that evaluation as open.
\end{remark}

\section{Conservation}
\label{sec:ch12-cons}

A constant \(\Lambda_{\mathrm{eff}}\) produced by
Theorem~\ref{thm:Leff} satisfies
\begin{equation}
  \label{eq:Leff-cons}
  \nabla^\mu(\Lambda_{\mathrm{eff}}g_{\mu\nu})=0
\end{equation}
automatically. The design constraint of
Section~\ref{sec:ch11-vacuum} is therefore met in the vacuum.

If the boundary data \(\phi\) are dynamical, \(S_{\mathrm{ent}}\)
can depend on time and \(\Lambda_{\mathrm{eff}}\) can drift. In
that case Theorem~\ref{thm:conservation} requires a compensating
current. The compensation is the boundary analogue of
Exercise~5 in Chapter~11: the variation of \(S_{\mathrm{HBEV}}\)
with respect to \(\phi\) supplies a boundary flux that, by the
on-shell identity for the radial mode, matches
\(\partial_\mu\Lambda_{\mathrm{eff}}\). The book uses a constant
\(\Lambda_{\mathrm{eff}}\) in all later solutions and models and treats
a rolling case as a cosmological refinement, not as the default.

\section{Relation to the diamond entropy of Chapter~6}
\label{sec:ch12-diamonds}

The two entanglement functionals in the book now have distinct
jobs.

\begin{center}
\begin{tabular}{@{}lll@{}}
\toprule
Functional & Support & Job \\
\midrule
\(S_D\) & small bulk diamonds & kill \(m_T\) \\
\(S_{\mathrm{HBEV}}\) & asymptotic boundary & set \(\Lambda_{\mathrm{eff}}\) \\
\(I(A:B)\) & two compact cores & conservative potential;
                                  phenomenological drag model not
                                  yet derived \\
\bottomrule
\end{tabular}
\end{center}

They are not interchangeable. Using \(S_D\) to subtract
\(\rho_\Lambda\) would confuse a local first-law identity with a
global vacuum-energy problem. Using \(S_{\mathrm{HBEV}}\) to kill
\(m_T\) would confuse a boundary condition with a bulk Ward
identity. The architecture is intentional.

\section{Working value used later}
\label{sec:ch12-working}

Later chapters use the regular-core model
\begin{equation}
  \label{eq:fr-with-L}
  f(r)
  =1-\frac{2M r^2}{r^3+r_0^3}
   -\frac{\Lambda_{\mathrm{eff}}}{3}r^2
\end{equation}
with \(\Lambda_{\mathrm{eff}}\) a constant of either sign. The
negative sign in \(f\) for \(\Lambda_{\mathrm{eff}}>0\) is
de Sitter-like and is the default cosmological reading; the
opposite sign is a mathematical option for matching interiors. No
Solar-System observable in Chapters~13--15 depends on the numerical
size of \(\Lambda_{\mathrm{eff}}\) at laboratory precision. In this
model, \(f(r)\to1-\Lambda_{\mathrm{eff}}r^2/3\) as \(r\to0\).

\section{Exercises}
\label{sec:ch12-exercises}

\begin{enumerate}
  \item \textbf{GHY variation.}
    Vary~\eqref{eq:GHY-again} with respect to \(h_{ij}\) and
    recover the \(K^{ij}-Kh^{ij}\) piece of~\eqref{eq:BY}.
  \item \textbf{Reference subtraction.}
    On Minkowski space in a large sphere of radius \(r_B\),
    compute \(K\) and \(\sqrt{h}\). What \(T_{\mathrm{bdy}}\)
    cancels \(S_{\mathrm{HBEV}}\) at \(\phi=0\),
    \(S_{\mathrm{ent}}=\alpha_0\LUV^3\)?
  \item \textbf{Why not \(\LUV^4\).}
    Contrast a four-dimensional counterterm
    \(\int_M\sqrt{-g}\,\LUV^4\) with a three-dimensional
    counterterm \(\int_{\partial M}\sqrt{h}\,\LUV^3\). Which
    tuning is a cancellation of two bulk quartics, and which is
    a renormalization of a boundary density?
  \item \textbf{Trace condition.}
    Reduce~\eqref{eq:bdy-EE} on a spherically symmetric tube
    \(r=r_B\) in Schwarzschild--de Sitter and show that it
    constrains \(\Lambda_{\mathrm{eff}}\) rather than \(M\).
  \item \textbf{Rolling \(\Lambda\).}
    Let \(S_{\mathrm{ent}}\) depend on a homogeneous boundary
    scalar \(\phi(t)\). Write the extra flux needed to keep
    Theorem~\ref{thm:conservation}. Compare with a slowly
    rolling quintessence potential.
  \item \textbf{Two functionals.}
    In one paragraph, explain what goes wrong if one attempts to
    set \(\Lambda_{\mathrm{eff}}\) from \(\delta S_D=0\) on a
    single small diamond.
\end{enumerate}

\section{Bridge to Chapter~13}
\label{sec:ch12-bridge}

The conditional infrared Einstein equation has the form
\begin{equation}
  \label{eq:EFE-ready}
  G_{\mu\nu}+\Lambda_{\mathrm{eff}}g_{\mu\nu}
  =8\pi\Geff\,T_{\mu\nu}^{\mathrm{total}},
\end{equation}
with \(G_{\mathrm{eff}}\) to be matched by the specified
operator, weight, threshold, continuation, and scheme calculation
of Problem~F, \(\Lambda_{\mathrm{eff}}\) from
Theorem~\ref{thm:Leff},
\(m_T=0\) from Theorem~\ref{thm:mT}, and conservation from
Theorem~\ref{thm:conservation}. What remains of the geometric
programme is the high-curvature interior, where the Einstein
truncation fails and Assumption~\ref{ass:ECG} is used. Chapter~13
uses the static spherical sector of Einsteinian cubic gravity
coupled to~\eqref{eq:EFE-ready} to motivate the regular-core model
\eqref{eq:fr-with-L}, which is a model and not an established
solution (Section~\ref{sec:status-J}), and only then turns to galactic
backreaction as a consequence of \(\Tent_{\mu\nu}\) in the weak-
field limit---not as a substitute for the derivation just
completed.